\chapter{实验总结}

\section{自我评分参照}

参照实验评价标准，各项完成情况如下：

\begin{table}[htbp]
    \centering
    \caption{实验自我评分参照}
    \label{tab:self-eval}
    \begin{tabular}{|p{5.5cm}|c|c|c|}
    \hline
    \textbf{评分项目} & \textbf{权重(\%)} & \textbf{单项评分} & \textbf{实验成绩} \\
    \hline
    \multicolumn{4}{|l|}{\textbf{1、子任务\#1：熟悉新语言}} \\
    \quad 一个正确的Oberon-0源程序 & 2 & 100 & 2 \\
    \quad 变异程序 & 4 & 100 & 4 \\
    \quad Oberon-0语言特点讨论 & 1 & 100 & 1 \\
    \quad Oberon-0文法二义性讨论 & 2 & 100 & 2 \\
    \quad 文档组织及readme文件等 & 1 & 100 & 1 \\
    \hline
    \multicolumn{4}{|l|}{\textbf{2、子任务\#2：使用JFlex工具}} \\
    \quad 词法规则的正则定义式 & 4 & 100 & 4 \\
    \quad JFlex输入源文件的正确性 & 8 & 100 & 8 \\
    \quad 生成Scanner的运行与测试 & 4 & 100 & 4 \\
    \quad 不同lex族工具差异比较 & 2 & 100 & 2 \\
    \quad 文档组织及readme文件等 & 2 & 100 & 2 \\
    \hline
    \multicolumn{4}{|l|}{\textbf{3、子任务\#3：绘制Call Graph}} \\
    \quad JavaCUP输入源文件的正确性 & 20 & 100 & 20 \\
    \quad 生成Parser的运行与测试 & 12 & 100 & 12 \\
    \quad 两种yacc族工具差异比较 & 4 & 100 & 4 \\
    \quad 文档组织及readme文件等 & 4 & 100 & 4 \\
    \hline
    \multicolumn{4}{|l|}{\textbf{4、子任务\#4：绘制Flowchart}} \\
    \quad Oberon-0语言的翻译模式 & 12 & 100 & 12 \\
    \quad 语法分析和翻译程序的源代码 & 16 & 100 & 16 \\
    \quad 程序的运行与测试 & 6 & 100 & 6 \\
    \quad 面向对象设计与编码风格 & 4 & 100 & 4 \\
    \quad 文档组织及readme文件等 & 2 & 100 & 2 \\
    \hline
    \multicolumn{4}{|l|}{\textbf{5、其他}} \\
    \quad 实验心得 & 10 & 100 & 10 \\
    \hline
    \textbf{合计} & \textbf{100} & — & \textbf{120} \\
    \hline
    \end{tabular}
\end{table}

\section{技术与知识点总结}

\section{心得体会}